\section{Introduction: what positivity remembers}
\label{sec:introduction}

A probability measure on pure quantum states has several mathematical lives.
It is an ensemble for an experimentalist, a positive integration functional
for a geometer, and a sequence of symmetric tensor moments for an operator
theorist. Passing between these descriptions is useful only if the relevant
structure survives. The support of the measure, the Hermitian metric on the
ambient vector space, and the positivity of the integration weights all matter.
For example, two ensembles can have identical two-copy statistics and very
different algebraic supports. A third-copy observable can then detect a
geometric distinction that is invisible to every two-copy measurement.

The same discipline is useful for many-body operators. A small transfer-map
error in an ordinary norm need not remain small after contraction with an
arbitrary tensor environment. Completely positive order does survive that
contraction. Likewise, a scalar cluster expansion need not control an ordered
operator logarithm, whereas a positive average of modular-time conjugations
preserves complete positivity term by term. The central theme of this paper is
to identify representations in which the required positivity is preserved and
then extract quantitative consequences from them.

\subsection{The problem and established approaches}

The motivating physical question asks when a local frustration-free
Hamiltonian has a spectral gap bounded away from zero as the system grows.
The AKLT chain has a classical rigorous solution through its valence-bond
structure \citep{q:AKLT1988}. In higher dimensions, decorated AKLT models
admit rigorous tensor-network gap proofs
\citep{q:ARLLNY2020,LuciaYoung2023}; boundary-state factorization provides
another route, developed for the hexagonal model in the recent preprint
\citet{Kastoryano2026}. These arguments make transfer operators, overlaps
of local ground spaces, and the normalization of the physical energy
central objects. We seek estimates that remain valid for arbitrary
positive tensor environments, so that a finite channel calculation can
certify a whole family of Hamiltonians.

The geometric question has an analogous emphasis on preserved structure:
what can a positive moment constraint force about support, concentration,
and equilibrium? Classical projective designs fix low tensor moments,
geometric invariant theory controls balanced tensors, and pluripotential
theory converts concentration into a nonlinear energy. Each has its own
established proof methods. The aim here is to build explicit interfaces
between them, retaining the Hermitian metric and positivity throughout.
This produces exact support classifications and quantitative certificates,
as well as an analytic entropy threshold.

For the associated particle systems, positive-temperature determinant
large deviations are established by \citet{prob_berman2017}. On complex
curves, a Vandermonde factorization also relates attractive determinant
ensembles to logarithmic pair interactions. In higher dimension the
determinant has collective singularities, and its microscopic pressure is
a distinct issue from macroscopic coercivity
\citep{AndreassonBermanSvensson2026}. We state the sign, constraint, and
scope of each probability result explicitly. The same care distinguishes
known scalar subordination laws from the positive operator approximations
and observable-dependent error bounds constructed from them.

\subsection{Three interfaces and eight principal contributions}

The first interface is between \emph{moments and algebraic support}.
The $t$th moment of an ensemble is a positive matrix on a symmetric tensor
power. A polynomial vanishing on the support gives a kernel vector of that
matrix. A Hermitian polynomial nonnegative on the support gives a separating
functional. These are concrete connections to geometric quantization,
Hermitian sums of squares, and quantum de Finetti theory
\citep{Rao2024,GouveiaParriloThomas2010}.
The second interface is between \emph{determinants and concentration}:
vanishing orders of evaluation determinants measure energy gained by
concentrating a particle ensemble. The third is between \emph{completely positive
order and local quantum constraints}: finite matrix inequalities become
certificates for large tensor networks without introducing the condition
number of the tensor core.

The contributions can be read in the following order.
\begin{enumerate}[label=\arabic*.]
\item \textbf{A sharp balanced entropy--energy coefficient.}
For first-moment-balanced measures on $\Pp^d$, the strict-subcritical
coercivity threshold is
\[
\Gamma_d^{\mathrm{bal}}=\frac{d}{d+1-H_{d+1}},
\qquad H_m=\sum_{j=1}^m\frac1j.
\]
The proof uses minimal jet determinants, a moment-map argument, a conditioned
determinantal coupling, and a fixed-degree quantization estimate. Harmonic
numbers enter through the largest coordinate of a uniform simplex point.
Positive moment repair and exact energy comparison then give smooth
constrained equilibria on the projective plane and a full large-deviation
principle for repulsive determinant ensembles in shrinking design tubes.

\item \textbf{Coherent moments control determinant ensembles and geometric potentials.}
For a positive polarized compact complex manifold, we prove the full
attractive determinant Laplace principle below the exponential integrability
threshold of normalized plurisubharmonic potentials. On $\mathbb P^q$ with
anticanonical polarization the range is $0\le\gamma<1/(q+1)$, and every
fixed-order shrinking design tube is allowed. The proof combines compact
one-particle conditional densities, coarse-grained entropy, and a two-sided
small-coupling bound that imports mean-energy information from the known
repulsive sign. A separate curvature argument gives finite-particle
concentration. For $D=q+1$ and $M_1=N^{-1}\sum_iP_{z_i}$, a projective
Ward identity yields the exact susceptibility
\[
 \mathbb E\left\|M_1-I/D\right\|_{\rm HS}^2
 =\frac{D-1}{DN(1-\gamma)}.
\]
The identity holds whenever the unconditioned partition is finite and
$\gamma<1$; the LDP range above requires no such additional assumption.
The product mean energy has a positive Hessian expansion in squared
coherent moments of rank-two interpolation pencils. Its convexity
allows differentiation of the mean-energy limit and gives a logarithmic
interpolation formula for the inverse complex Monge--Amp\`ere operator.
The resulting potential estimator has an explicit sampling-variance
bound, uniform in the polynomial level.

\item \textbf{Cubic support is controlled by three-qutrit balance.}
We classify all plane cubics supporting a positive projective second design
in their given Hermitian embedding. A quartic rank witness completes a
two-factor-balanced determinantal representation to a fully balanced tensor.
The $E_6$ Cartan theory then forces a unitary Hesse form. Smooth members have
a unique supporting design; the singular case gives orthogonal triangles
with an explicit Fourier description.

\item \textbf{Algebraic rigidity becomes an operational certificate.}
A rank-two observable on three matched qutrit copies, together with the
one-copy moment defect, bounds transport distance to the Hesse ensemble.
The copy number is minimal in this access model. The geometry of a
complete intersection supplies a pointwise error bound, and the simplex
structure of the finite ensemble controls its weights.

\item \textbf{A finite design determines a K3 surface.}
A ten-point Singer design yields three disjoint $(-3)$-curves on a rational
surface. A cyclic triple cover and three contractions produce a K3 surface
with transcendental lattice $A_2(7)$. The calculation makes the relation
between moment constraints, special point configurations, and lattice
geometry explicit.

\item \textbf{Coherent Gram spectra interpolate between ordinary and
noncrossing partitions.}
At unit sampling density in complex dimension $d$, the limiting $p$th moment is
\[
m_p(d)=\sum_{\pi\in\mathcal P(p)}\tau(G_\pi)^{-d}.
\]
Here $G_\pi$ is the directed cycle quotient associated with a set partition
and $\tau$ counts rooted directed spanning trees. Crossing partitions
survive at every fixed dimension, but disappear in the limit $d\to\infty$.
For logarithmic energy, an independent Newton-polytope argument replaces
torus averaging by a finite assignment with Kullback--Leibler cost, with
an explicit error of order $\log N/k$. This connects coherent determinants
to tropical geometry and optimal transport through a computation that
admits exact rational dual certificates.

\item \textbf{Positive operator calculus gives quantitative compilers.}
We relate ordered logarithms to order polynomials, represent compatible
modular responses by positive random times, and give local-commutator error
bounds for finite clock approximations. A separate completely positive
wire comparison controls angles between tensor-network constraints and
yields logarithmic sufficient AKLT decoration in the maximum graph degree.

\item \textbf{Moment contact controls measurement simulation.}
The same simplex expected maximum fixes the Haar-parent visibility
$\eta_D=(H_D-1)/(D-1)$. A quantitative anisotropic contact analysis gives
an open simulation neighborhood of projective-simulable measurements in
every fixed dimension and outcome number, and of the larger convex hull
of reducible qutrit measurements at visibility $5/12$.
\end{enumerate}

These results share tools, but their proof dependencies are deliberately
explicit. For instance, the spectral-gap certificate uses completely positive
comparison and subspace geometry. It does not require a thermodynamic limit
for a determinant ensemble. The Hesse certificate uses finite-dimensional
algebra and transport, so it can be implemented without solving a nonlinear
equilibrium problem.

\subsection{Position within the literature}

The use of symmetry breaking by moment constraints in Fano geometry is
developed by Andreasson--Berman--Svensson \citep{AndreassonBermanSvensson2026}.
Our balanced energy calculation concerns their analytic, orbit-reduced
coercivity framework. The determinant ensemble has additional microscopic
structure. Theorems~\ref{att:full-laplace}--\ref{att:polarized-extension}
give an attractive large-deviation principle below the section integrability
threshold; extension to the canonical coupling is posed in the final
section. This separates a geometric stability coefficient from a statement
about every singular stratum of a many-particle determinant.

The mean-energy approach to singular Gibbs measures, including its
positive- and negative-temperature distinction, is developed by
\citet{Berman2018meanenergy}. Our argument supplies the requisite entropy
and energy convergence for the actual attractive tilted determinant laws
in a definite range. In particular, it uses the established repulsive
theorem without assuming product-law energy convergence as an input.
For these Slater Hamiltonians, this supplies the smooth-reference
mean-energy limit considered in \citet[Problem 4.10]{prob_berman2017}.
The nonlinear Green formula in Theorem~\ref{cy:green-formula} is the
representation formulated in \citet[Conjecture 6.1]{prob_berman2017}.
A later conditional proof uses a complex zero-free neighborhood for a
two-parameter partition function
\citep[Theorems 3.2--3.3]{Berman2023zeta}. Our proof obtains the formula
from real mean-energy convergence and the positive projective-kernel
Hessian, without asserting that zero-free property. This is a concrete
connection between coherent tensor moments, interpolation theory, and
the nonlinear equation solved by Yau's theorem \citep{Yau1978}.
Toric pluripotential duality and complex optimal transport are established
in \citet{BermanWittNystrom2014}; the finite assignment error bound makes
that interface computational. The curvature estimate uses classical
H\"ormander theory, while the projective change of variables evaluates a
specific covariance exactly.

Moment-improved Sobolev and Moser--Trudinger inequalities have a substantial
history. Recent sharp results on spheres under finite-order moment constraints
are given by Zhou--Zou \citep{ZhouZou2026}. The energy here is the nonlinear
dual Monge--Amp\`ere energy, and the proof uses projective jets and balanced
determinantal couplings. The role of positive cubature remains a useful
conceptual comparison, with the geometric energy and its normalization kept
distinct.

The invariant theory used in the cubic classification is established:
Kempf--Ness converts zero moment into closed orbits and unitary uniqueness,
and the three-qutrit Cartan representation is treated explicitly by
Di Trani--de Graaf--Marrani \citep{geo_ditrani2023}. The additional bridge is
the rank witness that forces a design-supported pencil into that representation's
balanced locus. The Hesse pencil and its classical incidence geometry are
described by Artebani--Dolgachev \citep{geo_hesse2009}.

For quantum spin systems, tensor-network methods already prove gaps for
decorated AKLT models. Lucia--Young obtain a sufficient decoration linear
in the maximum degree \citep{LuciaYoung2023}. The quantitative comparison
developed here uses the positive order of the half-wire Choi matrices;
the cancellation between two ports gives logarithmic sufficient decoration.
The elementary finite-dimensional certificate and the transfer back to the
specified microscopic Hamiltonian are stated separately, including the
finite-patch gap factor.

The scalar laws underlying modular clocks are also established. Hyperbolic
infinite divisibility and Brownian subordination are treated by Pitman--Yor
\citep{PitmanYor2003}. We use them as positive operator representations and
combine them with observable-specific nested-commutator estimates. This
distinguishes the scalar sampling identity from the quantitative operator
approximation it supports.

\subsection{Reading routes and conventions}

Readers coming from quantum information can start with the moment operators
below, then read the cubic classification and ensemble certificate.
Readers from complex geometry can begin with the balanced energy theorem and
the Singer surface construction, then the nonlinear Green formula in
Section~\ref{cy:section}. The spectral-law and finite assignment
sections are largely self-contained after the definition of coherent
vectors. For the probability results, Section~\ref{att:entropy-section}
gives the entropy mechanism and the full attractive LDP, while
Section~\ref{ac:section} gives finite-particle bounds and the susceptibility
identity. The positive
calculus and network sections can be read independently using ordinary
finite-dimensional operator theory.

All logarithms are natural. Pure states are rays represented by unit vectors,
and $P_z=zz^*$ is independent of their phase. Our pure-state distance is
\[
d_{\mathrm{tr}}(P_z,P_w)=\tfrac12\|P_z-P_w\|_1
=\sqrt{1-|\langle z,w\rangle|^2}.
\]
On $\Pp^d$, Fubini--Study measure has total mass one, and the geometric
energy uses the anticanonical polarization $(d+1)H$. Spectral gaps of blocked
parents refer to sums of norm-one projections. The conversion to a
nearest-neighbor Hamiltonian keeps its local energy scale explicitly.

The supplemental source provides full proof modules for long geometric
arguments, exact arithmetic certificates, figure-generation code, and an
index of theorem dependencies. The verification section describes the
mathematical meaning of those computations and their trusted arithmetic
assumptions.
